% Appendix O: Civilization Dynamics and Collective Memory

\section*{O.1 Introduction}

Individuals learn. Civilizations remember. The central question is:
\emph{How does accessibility persist beyond individual lifetimes?}
A civilization is not merely a collection of people. It is a mechanism
for preserving trajectories.

\begin{definition}{Civilizational State}{civ-state}
  $C = (P, K, I, T)$ where $P$ is population, $K$ is knowledge,
  $I$ is infrastructure, and $T$ is transmission capacity.
\end{definition}

\begin{definition}{Collective Memory}{collective-mem}
  \emph{Collective memory} is the persistent storage of semantic structures
  beyond individual agents (books, libraries, archives, software repositories,
  institutions, traditions). The \emph{memory density field} $M_c(x)$
  captures its geographic distribution.
\end{definition}

\begin{definition}{Infrastructure as Accessibility Storage}{infra-access}
  Roads preserve transportation futures; power grids preserve energetic
  futures; libraries preserve semantic futures; universities preserve
  educational futures. \emph{Infrastructure is any structure preserving
  future accessibility.}
\end{definition}

\begin{definition}{Cultural Soliton}{cultural-soliton}
  A \emph{cultural soliton} is a stable localized semantic structure
  persisting for centuries (major religions, mathematical traditions,
  legal systems, foundational myths). They survive perturbation and
  propagate through populations.
\end{definition}

\begin{definition}{Civilizational Accessibility}{civ-access}
  $S_C = \log |\mathcal{A}_C|$, where $\mathcal{A}_C$ is the set of
  coherent futures available to the civilization. This becomes the
  central macro-observable.
\end{definition}

\begin{definition}{Accessibility Collapse}{acc-collapse}
  \emph{Accessibility collapse} occurs when $\frac{dS}{dt} \ll 0$:
  institutional breakdown, ecological failure, information loss,
  infrastructure decay. Collapse is not merely destruction —
  it is future contraction.
\end{definition}

\begin{theorem}{Adaptive Capacity Theorem}{adaptive-cap}
  Civilizations with larger accessible future volumes possess greater
  adaptive capacity.
\end{theorem}
\begin{proof}
  Larger accessibility implies more admissible responses to perturbation;
  more responses imply greater resilience.
\end{proof}

\begin{theorem}{Civilizational Persistence Theorem}{civ-persist}
  The long-term persistence of a civilization depends upon its ability
  to preserve and transmit accessibility across generations.
\end{theorem}
\begin{proof}
  Future generations inherit possibilities through transmission structures;
  loss of transmission reduces accessible futures; sustained transmission
  preserves future opportunity.
\end{proof}

\begin{namedthm}{The Accessibility Civilization Principle}
  The fundamental function of intelligence, education, science,
  infrastructure, and civilization is the preservation and expansion
  of coherent future accessibility:
  \[
    \text{Constraint} \to \text{Projection} \to \text{Meaning}
    \to \text{Knowledge} \to \text{Civilization} \to \text{Accessibility.}
  \]
  At every scale — from gestures to concepts, from minds to institutions,
  from semantic fields to cosmological fields — the same structure appears:
  \begin{center}
    \textbf{Persistent systems are accessibility-preserving systems.}
  \end{center}
\end{namedthm}
